\section{Limitations and Broader Impact}
\label{sec:limitations}

\paragraph{Statistical power and trigger scarcity.}
The evidence-rich set has 40 tasks.
Under the U-Net proposer, 720 reinspection episodes and 480 navigation
episodes all finish with zero strict success and zero reinspection triggers.
The controller is therefore never exercised, and the study cannot claim that
reinspection improves decisions or navigation.
A larger benchmark alone is insufficient; it must contain enough
evidence-positive episodes to exercise the policy.
This run also lacked a reachable VLM/LLM endpoint, so hybrid grounding
degraded to detector-only proposals.

\paragraph{Remaining robustness suites.}
GPS-noise and forced-degraded observation suites were not regenerated under
the U-Net proposer and are not reported as completed results.

\paragraph{Simulation and sensing fidelity.}
The environment uses georeferenced overhead imagery and approximates altitude changes by changing crop footprint.
It does not model wind, oblique perspective, motion blur, communications, battery dynamics, no-fly zones, or the full sensor pipeline of a physical UAV.
Performance in this environment is not evidence of safe autonomous deployment.

\paragraph{Perception bottlenecks.}
Damage labels can be ambiguous even to humans, and domain shift across satellite, crewed-aircraft, and low-altitude UAV imagery can invalidate confidence estimates.
The event-disjoint U-Net proposer mitigates the previous YOLO recall bottleneck for buildings, but damage grading, open-vocabulary grounding, and trigger coverage can still fail even when proposals are plentiful.
Calibration measured on one distribution does not guarantee calibration after shift.
Open-vocabulary vision-language outputs may improve coverage while introducing hallucinated evidence.

\paragraph{Responsible use.}
\ProjectName is intended as decision support.
False negatives can delay aid, while false positives can waste scarce response resources.
Human operators should retain authority over routing and damage declarations.
Geospatial imagery may expose sensitive locations; data access, retention, and released trajectories require disaster-specific privacy and security review.
